\uuid{UM47}
\titre{Distance}
\auteur{bodin ; hanani ; mzari}
\organisation{exo7}
\licence{CC BY-NC-SA 4.0}
\datecreate{2020-01-20}
\niveau{L1}
\module{Géométrie}
\chapitre{Géométrie de l'espace}
\sousChapitre{Distance}
\difficulte{1}
\theme{Distance}
\exo7qcmid{lille-1.geometrie-espace.023}

\contenu{
\texte{
Pour ces questions, l'espace est muni d'un repère orthonormé direct  $(O, \vec{i}, \vec{j}, \vec{k})$.
}

\begin{question}
\qclassification{141.05}
\qtags{motcle=distance}
Soit $A(2,0,1)$ et $D$ la droite d'équations :
$$\left\{\begin{array}{ccl}x+y-z&=&1\\x-y&=&-1 \end{array}\right.$$
  La distance de $A$ à $D$ est :
\begin{answers}

    \good {$\frac{3}{\sqrt 2}$}

    \bad {$\frac{1}{\sqrt 2}$}

     \bad{$\sqrt 3$}

    \bad{$\sqrt 2$}

\end{answers}
\begin{explanations}
Si $D$ est une droite qui passe par un point $B$ et dirigée par un vecteur $\vec{u}$  et $A$ un point, la distance de $A$ à $D$ est donnée par : $\frac{ \Vert \overrightarrow{BA} \wedge \vec{u} \Vert}{\Vert \vec{u}  \Vert}$.
\end{explanations}
\end{question}
}
